\subsection{Définition}
\begin{definition}
    Deux triangles sont \textbf{semblables} si les longueurs de leurs côtés sont proportionnelles.
\end{definition}
\begin{exemple*1}
    \titreExemple{Avec un quadrillage}

    \begin{minipage}{0.55\linewidth}        
        \begin{tikzpicture}[scale=0.45]
            \draw[help lines, color=black!30] (-5,-6) grid (10,0);                            
            \coordinate (T) at (-3,-5);
            \coordinate (W) at (5,-5);
            \coordinate (O) at (-4,-3);
            \coordinate (E) at (9,-3);
            \coordinate (P) at (-2,-1);
            \coordinate (B) at (1,-1);
            \tkzDrawSegment[thick,color=blue](O,T);
            \tkzDrawSegment[thick,color=blue](W,E);
            \tkzDrawSegment[thick,color=red](O,P);
            \tkzDrawSegment[thick,color=red](B,W);
            \tkzDrawSegment(T,P);
            \tkzDrawSegment(B,E);
            \tkzLabelPoints(T,W);
            \tkzLabelPoints[above](P,B,E);
            \tkzLabelPoints[left](O);
        \end{tikzpicture}
    \end{minipage}
    \hfill
    \begin{minipage}{0.45\linewidth}
        Sur la figure ci-contre, à l'aide du quadrillage, on peut constater que :
        \begin{itemize}
            \item $TO =2\times WE$
            \item $OP =2\times EB$
            \item $PT =2\times BW$
        \end{itemize}
        D'après la définition, puisque les longueurs de leurs côtés sont proportionnelles, \textbf{les triangles $\mathbf{TOP}$ et $\mathbf{WEB}$ sont donc semblables}.
    \end{minipage}
\end{exemple*1}

\begin{exemple*1}
    \titreExemple{Avec un tableau de proportionnalité}    

    \medskip
    \begin{minipage}{0.45\linewidth}
        Les triangles $MER$ et $GIF$ sont semblables. Le tableau suivant est donc un tableau de proportionnalité :

        \medskip
        \begin{tabular}{|c|c|c|c|}
        \hline 
        Triangle $MER$ & $ME$ & $ER$ & $RM$ \\ 
        \hline 
        Triangle $GIF$ & $GI$ & $IF$ & $FG$ \\ 
        \hline 
        \end{tabular}

        \medskip
        Les trois rapports suivants sont donc égaux : $\dfrac{ME}{GI}=\dfrac{ER}{IF}=\dfrac{RM}{FG}$
    \end{minipage}
    \hfill
    \begin{minipage}{0.5\linewidth}
        \begin{tikzpicture}[scale=0.5]
            \coordinate (M) at (0,4);            
            \coordinate (E) at (-2,2);
            \coordinate (R) at (-2,-1);
            \coordinate (F) at (-3,-3);
            \coordinate (I) at (6,-3);
            \coordinate (G) at (12,3);
            \tkzDrawSegment[color=red](E,M);
            \tkzDrawSegment[color=red](I,G);
            \tkzDrawSegment[color=blue](E,R);
            \tkzDrawSegment[color=blue](F,I);
            \tkzDrawSegment(R,M);
            \tkzDrawSegment(F,G);
            \tkzLabelPoints(R,F,I);
            \tkzLabelPoints[above](M,G);
            \tkzLabelPoints[left](E);
        \end{tikzpicture}
    \end{minipage}
\end{exemple*1}

\begin{remarque}
    \begin{itemize}
        \item Si deux triangles sont \textbf{homothétiques} alors ils sont \textbf{semblables}.
        \item La réciproque : \og Si deux triangles sont \textbf{semblables} alors ils sont  \textbf{homothétiques} \fg 
        
        est FAUSSE.
    \end{itemize}
\end{remarque}
